% ============================================================
% Week 04 section --- NLP for Structured HAZOP Reasoning
%
% Experimental investigation E1-E7 on the reconstructed T-301
% track. Metrics are against reconstructed row-level silver
% associations, not expert-adjudicated causal ground truth.
% ============================================================

\chapter[Experimental Evaluation and Findings]
{Experimental Evaluation and Findings}

\section{Objective}

This stage moves from the research design established in Week 03 to a
controlled evaluation of the implemented HAZOP relationship-recovery
prototype. The experiments compare retrieval and graph-based reasoning,
investigate evidence-support and provenance constraints, examine scenario
split variability, and measure sensitivity to controlled relationship
evidence removal.

The evaluation is deliberately restricted to what the available data can
support. The T-301 evaluation track contains 40 reconstructed records. The
held-out targets are consequence values co-listed with causes in the same
worksheet row. These are treated as silver associations, not independently
verified causal relationships.

\begin{glance}[Experimental stage at a glance]
Seven experiments (E1--E7) were completed. Graph reasoning achieved the
strongest deterministic F1 on the fixed seed-42 split. Subsequent experiments
showed that stricter support and provenance filters can sharply reduce
coverage, that results vary across scenario splits, and that graph recovery
degrades as training relationships are randomly removed. No experiment
establishes expert-adjudicated causal correctness.
\end{glance}

\section{Dataset and Evaluation Protocol}

The reconstructed T-301 track contains 40 scenario records. A deterministic
scenario-level 80/20 split with seed 42 assigns 32 scenarios to training and
8 to testing. The fixed test set produces 19 cause-to-consequence queries.
All records belonging to one scenario remain in the same partition. Held-out
target labels and graph edges are removed from the retrieval repository.

For each held-out cause, the silver target set consists of the consequences
co-listed in its source worksheet row. This construction measures recovery
of worksheet associations. It does not prove that every cause is physically
responsible for every consequence in the row.

\section{E1: Controlled Strategy Comparison}

The fixed-split comparison evaluates sparse TF--IDF retrieval, graph
reasoning, a deterministic graph-plus-TF--IDF candidate union, prompt-only
LLM reasoning, and hybrid graph/vector/LLM reasoning. The LLM methods were
evaluated on a deterministic cap of 15 queries; their denominators are
therefore distinct from the 19-query deterministic methods.

\begin{table}[htbp]
\centering
\caption[E1 fixed-split strategy comparison]{E1 results on the reconstructed
T-301 track. Metrics are against silver row associations. LLM methods were
evaluated on 15 of 19 queries.}
\label{tab:w4-e1}
\small
\begin{tabularx}{\textwidth}{@{}Xrrrrr@{}}
\toprule
\textbf{Method} & \textbf{N} & \textbf{Precision} & \textbf{Recall} &
\textbf{F1} & \textbf{Coverage} \\
\midrule
Prompt-only LLM & 15 & 0.0000 & 0.0000 & 0.0000 & 0.0667 \\
TF--IDF retrieval & 19 & 0.2784 & 0.4909 & 0.3553 & 0.8947 \\
Graph reasoning & 19 & 0.6486 & 0.8727 & 0.7442 & 0.8421 \\
Graph + TF--IDF union & 19 & 0.4068 & 0.8727 & 0.5549 & 0.8947 \\
Hybrid graph/vector/LLM & 15 & 0.3704 & 0.6383 & 0.4688 & 0.8667 \\
\bottomrule
\end{tabularx}
\end{table}

Graph reasoning achieved the highest deterministic micro-F1 on this split.
The graph-plus-TF--IDF union increased coverage but introduced additional
false positives and reduced F1. The hybrid method did not improve on graph
reasoning in the matched 15-query comparison and had substantially higher
latency in this run. The prompt-only method almost always abstained.

\section{E2--E5: Evidence and Provenance Constraints}

E2 varied the minimum number of distinct training scenarios supporting a
complete typed cause--deviation--consequence path. Requiring support from
two scenarios reduced F1 from 0.7442 to 0.1846 and coverage from 0.8421 to
0.3158. Thresholds of three or more yielded no eligible proposals.

E3 counted support for the cause-to-deviation and deviation-to-consequence
edges independently. Raising the cause-edge threshold from one to two
reduced recall from 0.8727 to 0.3455 and coverage from 0.8421 to 0.3158.
Raising only the consequence-edge threshold had a smaller effect. These
results identify sparse cause-edge support as a bottleneck in this dataset;
they do not establish that repeated co-occurrence is physical confirmation.

E4 compared direct and intermediate multi-hop paths. No eligible
intermediate paths were available under the evaluated data and constraints.
The combined condition therefore behaved like the direct-path condition.
This is an evidence-availability limitation, not proof that longer causal
chains are intrinsically ineffective.

E5 tested provenance metadata and source-coherence constraints. Metadata
presence thresholds did not change the baseline result. The stricter
same-scenario and same-source-row filters both reduced coverage to 0.2105
and F1 to 0.1791, with precision falling to 0.5000. Documentary coherence
alone did not identify causal truth and over-filtered potentially useful
cross-record evidence.

\section{E6: Scenario-Split Robustness}

E6 repeated the scenario-level 80/20 holdout for ten seeds (42--51). Across
the runs, mean micro-F1 was 0.6881 (standard deviation 0.0479), with a
minimum of 0.5825 and maximum of 0.7442. Mean coverage was 0.8746
(standard deviation 0.1283). The ten runs are repeated random holdouts,
not ten-fold cross-validation: their training partitions overlap, and four
scenarios never appeared in a test partition.

The result demonstrates meaningful sensitivity to the scenario partition.
The repeated-holdout mean is descriptive and should not be interpreted as
an independent-sample confidence interval.

\section{E7: Controlled Relationship-Evidence Masking}

E7 held the seed-42 scenario split fixed and randomly masked 10--40\% of
the 181 typed training relationships, using ten masking seeds. For each
seed, masks were nested across rates. The test labels were left unchanged.
The experiment completed 40 conditions across three deterministic methods.

\begin{table}[htbp]
\centering
\caption[E7 graph reasoning under evidence masking]{Mean graph-reasoning
metrics across ten random masks at each rate. The 0\% row is the E1
fixed-split baseline, not an additional E7 masking condition.}
\label{tab:w4-e7}
\small
\begin{tabularx}{\textwidth}{@{}rrrrr@{}}
\toprule
\textbf{Masked} & \textbf{Precision} & \textbf{Recall} &
\textbf{F1} & \textbf{Coverage} \\
\midrule
0\% & 0.6486 & 0.8727 & 0.7442 & 0.8421 \\
10\% & 0.6454 & 0.8109 & 0.7182 & 0.8000 \\
20\% & 0.6439 & 0.7382 & 0.6872 & 0.7316 \\
30\% & 0.6396 & 0.6582 & 0.6469 & 0.6842 \\
40\% & 0.6476 & 0.5727 & 0.6030 & 0.6158 \\
\bottomrule
\end{tabularx}
\end{table}

As masking increased, graph-reasoning recall, F1, and coverage declined,
while precision remained comparatively stable. This suggests that evidence
loss primarily reduces the number of recoverable silver associations and
increases abstention. The graph-plus-TF--IDF union retained higher coverage
but had lower F1 than graph reasoning at every masking rate. TF--IDF was
the weakest method in this experiment.

The masking procedure is a controlled simulation of evidence deletion from
loader-generated typed relationships. It is not a model of naturally
occurring missingness in industrial HAZOP reports.

\section{Cross-Experiment Discussion}

The completed experiments support several bounded observations:

\begin{enumerate}
    \item On the fixed reconstructed T-301 split, graph reasoning was the
    strongest tested deterministic method by micro-F1.
    \item Candidate union increased coverage but traded away precision and
    F1; it is not a learned fusion or reranking method.
    \item Repeated support constraints reduced recovery because complete
    paths and cause edges were sparsely represented in the training data.
    \item Provenance-field completeness did not imply evidence strength,
    and strict same-source coherence removed useful proposals without
    improving precision.
    \item Intermediate multi-hop reasoning was not meaningfully exercised
    because no eligible intermediate paths were present.
    \item Repeated holdouts showed scenario-split sensitivity.
    \item Random relationship masking reduced graph-reasoner recall,
    F1, and coverage as evidence was removed.
\end{enumerate}

These findings concern agreement with reconstructed row-level associations.
They do not establish physical causality, expert approval, or safety utility.

\section{Limitations and Remaining Work}

The principal limitation is the absence of an independently
expert-adjudicated held-out relationship set. The reconstructed dataset
contains only 40 scenarios, and the test partition contains 8 scenarios
and 19 queries. The E6 repeated holdouts are not independent folds.
The E4 intermediate-path condition had no eligible examples. The E7 masking
experiment removes relationships uniformly at random and does not reproduce
real documentation missingness. LLM evaluation was capped at 15 queries,
and reliable provider cost data was not captured.

The next research requirement is independent process-safety adjudication
using a written rubric for causal plausibility, context sufficiency, and
safeguard validity, with reviewer disagreement and inter-rater agreement
reported. End-to-end validation of the desktop application and a complete
reproducibility audit are also required before final submission.

\begin{constraint}[Research and safety boundary]
The prototype is a research aid for structured retrieval and reasoning.
Its recovered links must not be treated as verified process-safety facts,
operational recommendations, or a substitute for qualified HAZOP review.
\end{constraint}

\section{Conclusion}

The E1--E7 investigation indicates that typed graph reasoning is a useful
representation and retrieval strategy for recovering row-derived HAZOP
associations in the reconstructed T-301 case. Its performance is sensitive
to the availability and distribution of training relationships. Stronger
support and source-coherence filters can sharply reduce coverage, while
random evidence masking reduces recall and increases abstention. These
results motivate expert-adjudicated validation and broader, better-grounded
datasets before stronger causal claims can be made.
